%% Evaluation protocol (MedIA revision, v2 pipeline).
%% Every implementation detail below is taken from src/v2/ (run_v2.py, arms.py,
%% select_hparams.py, stats_v2.py, analyze_v2.py); see the provenance list kept
%% with the revision notes.
\section{Evaluation protocol}
\label{sec:protocol}

A cohort has $K$ domains, and fold $t$ holds out domain $t$ for testing. Every
arm (a training method or a test-time procedure) is run in every fold with
several seeds. The cohorts, models and arms are described in
Section~\ref{sec:setup}.

\subsection{Leave-one-domain-out design}
\label{sec:protocol-lodo}

Leave-one-domain-out evaluation is standard in domain generalization
\citep{gulrajani2021domainbed}, and Camelyon17-WILDS adds a held-out validation
hospital \citep{koh2021wilds}. We follow both. On Camelyon17, the canine
cohort and MIDOG++, fold $t$ also holds out domain $(t-1)\bmod K$ as an
out-of-distribution (OOD) validation domain and trains on the remaining $K-2$
domains. MIDOG~2021 has three labeled scanners, so each fold trains on two and
has no OOD validation domain. Every domain is the test domain of exactly one
fold, so the evaluation is \emph{domain-complete}: no single test domain is
privileged. In every training domain, an in-distribution (ID) validation split
is held out that is disjoint from training in its groups (patients, slides or
images; Section~\ref{sec:setup-cohorts}). The study adds four elements to this
design.

\paragraph{Seed replication} Every arm is run in every fold with five seeds at
the compact tier and three at the standard tier
(Section~\ref{sec:setup-models}). The seed fixes the initialization of
randomly initialized layers, the batch order, the augmentation draws and the
order of the test stream for test-time adaptation. All arms use the same seed
values.

\paragraph{Fixed step budget} Every fine-tuned arm is trained for 1,200 steps
with batch size 128 and a cosine learning-rate schedule. It is evaluated at
ten checkpoints, at steps $\mathrm{round}(kn/10)$ for $k=1,\dots,10$ and a
budget of $n$ steps. The only exception is the compute-matched control
(Section~\ref{sec:setup-arms}), which needs a larger budget by design. With a
fixed budget, schedule length and early stopping cannot confound method
comparisons.

\paragraph{All model-selection rules} At each checkpoint we evaluate the ID
validation, OOD validation and test splits. Four rules select a checkpoint by
the maximum area under the ROC curve (AUROC) on their selection split, with ties
going to the earliest checkpoint:
\begin{enumerate}
  \item \emph{training-domain (ID) validation}, the primary rule
    \citep[DomainBed's default;][]{gulrajani2021domainbed}, available in every
    cohort;
  \item \emph{leave-one-domain-out validation} on the OOD validation domain
    (Camelyon17, canine, MIDOG++);
  \item \emph{last}, the final checkpoint;
  \item \emph{oracle}, the best test AUROC, reported only as an upper bound.
\end{enumerate}
All four rules use the same ten checkpoints of the same runs, so any
difference between them comes from selection alone.

\paragraph{Per-fold hyperparameter selection} Hyperparameters are tuned with
seed~0 on the grids of Table~\ref{tab:arms}. Every fold keeps the configuration
with the highest ID-validation AUROC at the ID-selected checkpoint of its own
tuning run. Diverged runs are excluded unless every configuration diverged, and
test data are never used. We do not average ID-validation AUROC over folds,
because doing so leaks test-domain labels. The test domain of fold $t$ is a
training domain of the other folds, except for fold $t+1$ when a validation
domain exists. The ID-validation splits of those folds therefore contain
labeled patches of domain $t$. In every cohort, some of these patches are
pixel-identical to rows of fold $t$'s test set.
Averaging would thus let fold $t$'s test labels influence the hyperparameters
reported for fold $t$. At the standard tier, the base learning rate is selected
first, and the method grids are then run at each fold's selected rate. The
Barlow Twins fine-tunes have their own rate grid (Section~\ref{sec:setup-optim}).

\subsection{Reported quantities}
\label{sec:protocol-metrics}

Let $y_{ts}$ be the test AUROC (the primary metric) of one arm in fold $t$
with seed $s=1,\dots,n_t$ under a given selection rule, and
$\bar y_t=n_t^{-1}\sum_s y_{ts}$. For every arm, cohort, tier and rule we report
\begin{align}
  \text{worst domain:}\quad & \min_t \bar y_t, \label{eq:worst}\\
  \text{across-domain SD:}\quad & s_{\bar y}=\Big[\tfrac{1}{K-1}\textstyle\sum_t(\bar y_t-\hat\mu)^2\Big]^{1/2},
     \label{eq:sdacross}\\
  \text{mean:}\quad & \hat\mu = \tfrac{1}{K}\textstyle\sum_t \bar y_t,\qquad
     \hat\mu \pm t_{0.975,\,K-1}\, s_{\bar y}/\sqrt{K}. \label{eq:mean}
\end{align}
With balanced seeds, the interval in \eqref{eq:mean} is exact under model
\eqref{eq:oneway}, since $\operatorname{Var}(\bar y_t)=\sigma_D^2+\sigma_R^2/n$
for every domain. Percentile bootstrap intervals \citep{efron1993bootstrap}
are supplementary only, because the domain bootstrap under-covers for $K\le7$.
Calibration is reported as the expected calibration error
\citep{naeini2015ece,guo2017calibration}, with 15 equal-width bins $B_b$ on the
confidence of the predicted class (decision threshold 0.5):
\begin{equation}
  \mathrm{ECE}=\sum_{b=1}^{15}\frac{|B_b|}{N_{\text{test}}}\,
  \big|\operatorname{acc}(B_b)-\operatorname{conf}(B_b)\big|,
  \label{eq:ece}
\end{equation}
summarized by its mean and maximum over domains.

\subsection{Variance decomposition}
\label{sec:protocol-variance}

For one arm on one cohort we fit the one-way random-effects model
\begin{equation}
  y_{ts}=\mu+b_t+e_{ts},\qquad b_t\sim\mathcal N(0,\sigma_D^2),\quad
  e_{ts}\sim\mathcal N(0,\sigma_R^2),\qquad s=1,\dots,n_t,
  \label{eq:oneway}
\end{equation}
with all terms independent. Seeds are nested within domains: seed $s$ in
folds $t$ and $t'$ denotes different runs on different training sets, so no
crossed seed effect is fitted. $\sigma_D$ is the between-domain and $\sigma_R$
the between-run SD. With $N=\sum_t n_t$, between- and within-domain mean
squares $\mathrm{MSB}=\sum_t n_t(\bar y_t-\bar y)^2/(K-1)$ and
$\mathrm{MSW}=\sum_{t,s}(y_{ts}-\bar y_t)^2/(N-K)$, and
$n_0=(N-\sum_t n_t^2/N)/(K-1)$ ($n_0=n$ for balanced data), the ANOVA
estimators \citep{searle1992variance} are
\begin{equation}
  \hat\sigma_R^2=\mathrm{MSW},\qquad
  \hat\sigma_D^2=\max\{(\mathrm{MSB}-\mathrm{MSW})/n_0,\,0\},\qquad
  \widehat{\mathrm{ICC}}=\frac{\mathrm{MSB}-\mathrm{MSW}}{\mathrm{MSB}+(n_0-1)\,\mathrm{MSW}},
  \label{eq:anova}
\end{equation}
where $\mathrm{ICC}=\sigma_D^2/(\sigma_D^2+\sigma_R^2)$ is ICC(1)
\citep{shrout1979icc,mcgraw1996icc}, reported truncated at 0. Restricted
maximum likelihood (REML) estimates \citep{patterson1971reml} are obtained by
profiling out $\sigma_R^2$ and minimizing over $\gamma=\sigma_D^2/\sigma_R^2\ge0$,
with the boundary $\gamma=0$ checked explicitly. They coincide with ANOVA for
balanced data with a positive estimate.

\paragraph{Intervals and boundary estimates} We report the exact $F$-based
95\% interval for the ICC, the exact $\chi^2$ interval for $\sigma_R$ and the
modified large-sample interval for $\sigma_D$
\citep{searle1992variance,mcgraw1996icc,burdick1992confidence}. With $K\le7$,
$\hat\sigma_D$ is often at or near zero, so we also report a Bayesian grid
posterior under weakly informative half-normal and half-Cauchy priors (scale
0.1 AUROC) \citep{gelman2006prior}, with one-sided intervals $[0,q_{0.95}]$ at
the boundary. The formulas are in Supplementary
Section~\ref{S-sup:protocol}.

\paragraph{The ICC as a diagnostic} The ICC is the share of an arm's
single-run variance that is due to which domain is held out. Near one, the
choice of domains dominates. Near zero, run-to-run noise is as large as the
differences between domains, and a single-seed result says little about a
domain. The ICC is conditional on arm, cohort, metric, selection rule and the
domains at hand, and with $K\le7$ its intervals are wide. It tells us where the
variance of a reported number comes from. It is not a quality standard, and we
set no threshold for it. As an external reference point we apply the same
estimators to DomainBed's published per-environment results
\citep{gulrajani2021domainbed}. The mean $\pm$ standard error over three
trials is converted to an SD with the convention of DomainBed's released code
(Supplementary Section~\ref{S-app:domainbed}).

\subsection{Paired contrasts and fold dependence}
\label{sec:protocol-contrasts}

Each arm is compared with ERM of the same tier and cohort through
$d_t=\bar y^{a}_t-\bar y^{\text{ERM}}_t$. We report $\bar d$, the interval
$\bar d\pm t_{0.975,K-1}s_d/\sqrt K$, the two-sided $p$-value and the number of
improved domains. The $p$-values are Holm-adjusted over all arms compared with
ERM within a cohort, tier and selection rule \citep{holm1979}.

The folds of a leave-one-domain-out design share training domains, so the
$d_t$ are not independent \citep{nadeau2003inference,bates2023cv}. Let $T_i$ be
the set of training domains of fold $i$ and
$O_{ij}=|T_i\cap T_j|/\sqrt{|T_i||T_j|}$ the overlap of folds $i$ and $j$. As
a sensitivity analysis we assume a working correlation
$\operatorname{corr}(d_i,d_j)=R_{ij}=\rho\,O_{ij}$ for $i\ne j$
($R_{ii}=1$). Under this assumption
$\operatorname{Var}(\bar d)=\sigma^2\,\mathbf 1^\top R\mathbf 1/K^2$, and the
sample variance is biased downward,
$\mathbb E[s_d^2]=\sigma^2(1-\bar r)$, where $\bar r$ is the mean off-diagonal
element of $R$. The corrected standard error and the effective number of folds
are then
\begin{equation}
  \widehat{\mathrm{SE}}_\rho(\bar d)=\Big[\frac{s_d^2}{1-\bar r}\,
  \frac{\mathbf 1^\top R\mathbf 1}{K^2}\Big]^{1/2},\qquad
  K_{\text{eff}}=\frac{K^2}{\sum_{i,j}R_{ij}},
  \label{eq:keff}
\end{equation}
with $t_{K-1}$ quantiles, in the spirit of the corrected resampled $t$
statistic of \citet{nadeau2003inference}. Because $\rho$ cannot be estimated
from $K$ folds, we report $\rho\in\{0,0.25,0.5,1\}$. For our designs,
$K_{\text{eff}}=3.33$, 2.50 and 1.67 at $\rho=0.25$, 0.5 and 1 for Camelyon17
and the canine cohort ($K=5$); 3.50, 2.33 and 1.40 for MIDOG++ ($K=7$); and
2.40, 2.00 and 1.50 for MIDOG~2021 ($K=3$).

\subsection{Power}
\label{sec:protocol-power}

Power for a paired comparison depends on the variance of the paired
difference, not on either arm's between-domain SD. Writing
$y^{a}_{ts}=\mu_a+b_t+(ab)_{t}+e^{a}_{ts}$ and
$y^{\text{ERM}}_{ts}=\mu_{\text{ERM}}+b_t+e^{\text{ERM}}_{ts}$, the domain
effect $b_t$ cancels, and with $s$ seeds per arm
\begin{equation}
  \sigma_\Delta^2(s)=\operatorname{Var}(d_t)=\sigma_{DA}^2+\sigma_{R,\Delta}^2/s,
  \qquad \sigma_{R,\Delta}^2=\operatorname{Var}(e^{a}_{ts}-e^{\text{ERM}}_{ts}),
  \label{eq:diffvar}
\end{equation}
where $\sigma_{DA}^2$ is the domain-by-arm interaction variance. Because all
arms share seed values, their run noise may be correlated, so
$\sigma^2_{R,\Delta}$ is estimated as the pooled within-domain variance of the
per-seed differences $y^{a}_{ts}-y^{\text{ERM}}_{ts}$, which is unbiased either
way, and $\hat\sigma_{DA}^2=\max\{s_d^2-\hat\sigma^2_{R,\Delta}h,0\}$ with
$h=K^{-1}\sum_t n_t^{-1}$. If seed sets are not aligned, we use
$\hat\sigma^2_{R,\Delta}=\mathrm{MSW}_a+\mathrm{MSW}_{\text{ERM}}$ and
$\hat\sigma_{DA}^2=\max\{s_d^2-\mathrm{MSW}_a h_a-\mathrm{MSW}_{\text{ERM}}h_{\text{ERM}},0\}$.
The power of the two-sided paired $t$ test over $K$ domains to detect a mean
difference $\delta$ is
\begin{equation}
  1-\beta(K)=\Pr\big(|T_{K-1}(\lambda)|>t_{1-\alpha/2,\,K-1}\big),\qquad
  \lambda=\frac{\delta\sqrt{K_{\text{eff}}}}{\sigma_\Delta(s)},
  \label{eq:power}
\end{equation}
where $T_{K-1}(\lambda)$ is noncentral $t$ with $K-1$ degrees of freedom and
$K_{\text{eff}}=K$ for independent folds. We report the smallest $K$ that
reaches 80\% power at $\alpha=0.05$, for
$\delta\in\{0.02,0.05,0.10,0.20\}$ AUROC and $s\in\{1,3,5,10\}$. For
$\rho>0$, $\sigma_\Delta$ is first divided by $\sqrt{1-\bar r_{\text{obs}}}$
(the bias correction of \eqref{eq:keff} for the observed design). The primary
\emph{ratio} mode then holds the observed inflation $K/K_{\text{eff}}$ fixed as
$K$ grows. The \emph{design} mode uses $K_{\text{eff}}$ of a hypothetical
$K$-domain design, which levels off near $1/\rho$ because training sets overlap
more as $K$ grows, so some targets are unreachable.

\subsection{Rank stability and the task-by-shift decomposition}
\label{sec:protocol-ranks}

Within a cohort, arms are ranked by worst-domain \eqref{eq:worst} and by mean
AUROC, with 95\% percentile rank intervals and the probability of ranking first
from 5,000 bootstrap replicates. Two resampling schemes are used: seeds within
each (arm, domain) cell, and domains with one shared draw for all arms.
Agreement between two cohorts' rankings is Kendall's $\tau_b$ over their common
arms \citep{kendall1938}, with a 95\% interval from 2,000 seed-bootstrap
replicates and a $p$-value against independent orderings. Of the six cohort
pairs, two share the task, two the shift type and two neither. We report the
mean worst-domain $\tau_b$ per class and bootstrap intervals of the differences, which are
descriptive.

\paragraph{Arm-by-cohort interaction} For the $A$ arms that are run in all
four cohorts, let $G_{cta}=\bar y^{a}_{ct}-\bar y^{\text{ERM}}_{ct}$ be the
gain of arm $a$ over ERM in domain $t$ of cohort $c$. Each (cohort, domain) row
is centered across arms, $E_{cta}=G_{cta}-A^{-1}\sum_{a'}G_{cta'}$, removing
the domain's average gain. Whether the centered profiles differ between cohorts
is tested with domains as replicates, using a studentized Welch--James statistic
and a wild bootstrap that multiplies whole residual rows by Rademacher signs
(10,000 replicates). This keeps each cohort's variance and the correlation
between arms.

\paragraph{Task, shift and task $\times$ shift} Task and shift type are
constant within a cohort, so they can only be assessed with the cohort as the
unit. Let $\mathbf m_c\in\mathbb R^A$ be the mean of the centered rows of
cohort $c$, with equal weight for every cohort, and let $s^{f}_c=\pm1$ code
cohort $c$ for factor
$f\in\{\text{task},\text{shift},\text{task}\times\text{shift}\}$, with
$s^{\text{task}\times\text{shift}}_c=s^{\text{task}}_c s^{\text{shift}}_c$.
Because the four cohorts form a $2\times2$ design with one cohort per cell, the
between-cohort interaction sum of squares splits exactly into three orthogonal
contrasts:
\begin{equation}
  \sum_{c}\|\mathbf m_c-\bar{\mathbf m}\|^2
  =\sum_{f}\mathrm{SS}_f,\qquad
  \mathrm{SS}_f=\|\mathbf L_f\|^2,\quad
  \mathbf L_f=\tfrac12\sum_c s^{f}_c\,\mathbf m_c .
  \label{eq:2x2}
\end{equation}
We report each $\mathrm{SS}_f$ and its fraction of the total, and the rank
of the task and shift splits among the three possible two-versus-two splits of
four cohorts. A cohort-level permutation $p$-value equals rank$/3$ and cannot
be below $1/3$. We also give an approximate test $F_f=\mathrm{SS}_f/\mathrm{SS}_{\text{task}\times\text{shift}}$
on $(A-1,A-1)$ degrees of freedom for $f\in\{\text{task},\text{shift}\}$. It
treats the cohort-level deviations as independent and spherical across arms.
Its error term has only $A-1$ degrees of freedom and absorbs any genuine
task-by-shift pattern. With four cohorts, this decomposition describes the
interaction. It cannot show that rankings are organized by task or by shift
type, and we do not claim that they are.

\subsection{A site acceptance test for new sites}
\label{sec:sat}

The diagnostics above describe a method across sites. A practitioner who is
about to use one trained model at a new hospital or scanner needs a decision
for that site: is the model's AUROC there at least a target $T$? We propose an
empirical-Bayes site acceptance test that combines the development evidence
with a small labeled sample from the new site.

\paragraph{Development prior} Let $z_{sr}=\operatorname{logit}(y_{sr})$ be the
test AUROC of run $r$ of the method at development site $s$ on the logit scale,
for $K'$ held-out development sites with $n$ runs each. Under model
\eqref{eq:oneway} on this scale, with a flat prior on $\mu$ and independent
half-Cauchy priors on $\sigma_D$ and $\sigma_R$, the AUROC of one run at a new
site has, on the logit scale, the posterior predictive distribution
\begin{equation}
  \theta_{\text{new}}\mid\sigma_D,\sigma_R,\mathbf z\ \sim\
  \mathcal N\!\Big(\bar z,\ \sigma_D^2+\sigma_R^2+\frac{\sigma_D^2+\sigma_R^2/n}{K'}\Big),
  \label{eq:sat-prior}
\end{equation}
integrated over the grid posterior of $(\sigma_D,\sigma_R)$, where $\bar z$ is
the grand mean. The result is a scale mixture of normal distributions centered
at $\bar z$. Its spread grows with the between-site SD and shrinks only slowly
with $K'$, so few or heterogeneous development sites give a vague prior. The
half-Cauchy scale is 0.5 on the logit scale (about 0.06 AUROC at an AUROC of
0.85).

\paragraph{Local evidence} At the new site, $m$ patches are labeled. Their AUROC
$a_m$ (with a continuity correction that keeps it inside $(0,1)$) gives
$\hat z_m=\operatorname{logit}(a_m)$ with variance
$v_m=\operatorname{Var}(a_m)/\{a_m(1-a_m)\}^2$, where $\operatorname{Var}(a_m)$
is the Hanley--McNeil variance for the observed numbers of positives and
negatives. Each mixture component is updated conjugately with
$\hat z_m\sim\mathcal N(\theta,v_m)$, and the component weights are updated by
their marginal likelihoods.

\paragraph{Decision} The site is accepted if
$P(\theta\ge\operatorname{logit}T\mid\mathbf z,\hat z_m)\ge 1-\alpha$, with
$\alpha=0.05$. The posterior also gives a point estimate and a credible interval
for the site's AUROC. Before any labels are collected, the number of labels
needed is planned by pre-posterior simulation: sites are drawn from the
development prior \eqref{eq:sat-prior}, local scores are simulated, and $m$ is
the smallest sample size for which sites that exceed the target by a chosen
margin (0.02) are accepted with the chosen probability (0.8). The simulation
uses equal-variance binormal scores at a prevalence of 0.5.

\paragraph{Validation design} Every held-out site of every cohort plays the
new site in turn, for every method and seed (the six extra frozen backbones of
the representation ladder excepted). Its prior is built only from the other
folds' results of the same method at their own held-out sites, so the new
site's test labels are never used. Samples that contain only one class are
discarded (at most 3 of 200). Labeled samples of size $m\in\{25,50,100,200,400\}$ are drawn at
random, 200 times, from the run's saved test predictions, and the truth is the
run's AUROC on the full test split. We compare the test with two baselines:
local labels alone (accept if the one-sided 95\% Wald bound of $\hat z_m$ exceeds
$\operatorname{logit}T$) and the development prior alone ($m=0$). We report
the rate of accepting a site whose true AUROC is below $T$ (false acceptance)
and the rate of accepting a site at or above $T$ (power), for
$T\in\{0.80,0.85,0.90\}$.
